On réalise l’électrolyse d’une solution de bromure de cuivre dans un
électrolyseur à électrodes de carbone.

La caractéristique $U=f(I)$ de cet électrolyseur peut être modélisée par le
graphique ci-contre.

On peut donc écrire~: $U=E+r I$.

\begin{figure}[H]
    \centering
    \begin{tikzpicture}
        \begin{axis}[
                cycle list name=my color list,
                width=9cm,height=6cm,
                xlabel=$I~(\unit{\mA})$,ylabel=$U~(\unit{\volt})$,
                xmin=0,xmax=245,
                ymin=0,ymax=3.5,
                xtick distance=50,
                ytick distance=1,
                domain=0:200
            ]
            \addplot {0.012*\x+0.8};
        \end{axis}
    \end{tikzpicture}
\end{figure}

\begin{enumerate}
    \item Déterminer la valeur de $E$ (préciser l'unité) et celle de la
          résistance interne $r$ de l'électrolyseur.
    \item On maintient aux bornes de l'électrolyseur une tension
          $U=\qty{2,0}{\volt}$ pendant une durée $\Delta t=20$ minutes.
          Quelle est l'énergie chimique fournie~?
    \item Quel est le rendement de l'électrolyseur~?
\end{enumerate}


